
\textbf{Corollaire 5.4.} \uline{La filtration induite par} @@M0@@ \uline{sur} @@M1@@ \uline{est identique à la filtration de} @@M2@@.

\textbf{Corollaire 5.5.} \uline{Soient} @@M3@@ \uline{un schéma quasi}-\uline{compact}, @@M4@@ \uline{des} @@M5@@-\uline{Modules} \uline{localement libres de type fini}, @@M6@@ \uline{le schéma de drapeaux de} @@M7@@, @@M8@@. \uline{Alors l'homomorphisme canonique} :
@@M9@@
\uline{est injectif}.

Le raisonnement de 3.1 permet de se ramener à montrer la même propriété quand @@M10@@ est un fibré projectif sur @@M11@@. Ce cas résulte alors immédiatement de 5.4.

\textbf{Corollaire 5.6.} \uline{Avec les notations et les hypothèses de} 5.1, @@M12@@ \uline{est un} @@M13@@-\uline{module libre}, \uline{admettant pour base} @@M14@@. \uline{Pour tout} @@M15@@, \uline{on a} :
@@M16@@
\uline{Enfin}, \uline{les éléments} @@M17@@ \uline{sont liés par la relation} :
@@M18@@
\uline{où} @@M19@@ \uline{est la} @@M20@@-\uline{ième classe de Chern de} @@M21@@ (V 6.7).

Conséquence immédiate de 1.1, 5.3, et (1.13.1).

\textbf{Corollaire 5.7.} \uline{Avec les notations et les hypothèses de} 4.5, \uline{la} @@M22@@-\uline{algèbre} @@M23@@ \uline{est engendrée par les} @@M24@@, \uline{où} @@M25@@ \uline{est la classe dans} @@M26@@ \uline{de} @@M27@@. \uline{L'idéal des relations entre les} @@M28@@ \uline{est engendré par les relations déduites de} :
@@M29@@
